\chapter*{Conclusion} % chapter* je necislovana kapitola
\addcontentsline{toc}{chapter}{Conclusion} % rucne pridanie do obsahu
\markboth{Conclusion}{Conclusion} % vyriesenie hlaviciek

This thesis has examined how machine-learning analysis of the textual record of \gls{ecb} press conferences quantifies, structures, and ultimately prices the information that the central bank communicates to financial markets. Working with a custom corpus of $289$ press conferences from June~1998 to October~2024, and a pipeline that separates the prepared \gls{is} from the spontaneous \gls{qa} and, within the \gls{qa}, journalists' questions from the President's answers, the empirical chapters delivered four findings.

First, FinBERT and CentralBankRoBERTa produce internally coherent series whose turning points align with the canonical episodes of the euro-area cycle: the October~2008 rate cut, Draghi's 2012 ``Whatever it takes'' speech, the 2015 \gls{qe} launch, the 2020 pandemic emergency, and the 2022 hiking lift-off. Two encoder families trained on different corpora agree on the sign of the signal, which is the minimum bar for treating the textual measure as more than a model artefact.

Second, the \gls{is} and the \gls{qa} are structurally different communicative acts. The \gls{is} traces a coherent three-act narrative arc---moderate opening, risk-catalogue trough, hawkish policy-decision close---with a within-conference swing of roughly $0.4$ FinBERT units, while the \gls{qa} is essentially flat as a function of relative position. Within the \gls{qa}, journalist questions are systematically more dovish than the \gls{ecb}'s answers, by an average of $0.29$ on the CentralBankRoBERTa scale, and this gap is stable across the Pre-\gls{zlb}, \gls{zlb}, and hiking regimes. Without separating the two voices, the press corps' affective tone would contaminate the policy signal by approximately $0.3$ units---an order of magnitude that is large relative to the policy-relevant variation. The structural decomposition of the press conference is therefore not an optional refinement but a methodological prerequisite for any sentiment-based study of central-bank communication, and is one of the principal departures of this thesis from earlier studies that treat the press conference as a single document.

Third, sentiment leads the \gls{ecb}'s policy decisions rather than tracking them coincidentally. The event-study window of $\pm 6$ meetings shows a hawkish profile fully established up to six meetings before an actual hike in the pre-\gls{zlb} period and a symmetric telegraphing of both hikes and cuts in the post-pandemic cycle. The lag-correlation analysis identifies a peak between the FinBERT overall mean and the \gls{mro} change at a lag of four meetings ($r = 0.551$ on changed meetings).

Fourth, the textual signal carries incremental predictive content for the direction of the next \gls{ecb} rate decision. Four classifiers were compared on a three-class target: a rate-history baseline (M1), a sentiment-only logit (M2), an augmented logit (M3), and a Random Forest on the same feature set (M4). The cross-validated AUC ordering is strict: $0.682 < 0.694 < 0.705 < 0.763$. The M2-versus-M1 increment establishes that the text adds directional information beyond the rate history alone; the M4-versus-M3 increment of $+0.058$ shows that an economically meaningful portion of that information is carried by non-linear interactions which the logistic specification cannot capture. The textual record and the rate history are complementary rather than redundant.

Four extensions are particularly natural. First, the affective-sentiment layer should be replaced with a dedicated \emph{stance classifier} that assigns the operationally meaningful hawkish/dovish/neutral label rather than positive/negative tone; the \gls{zlb}-era flatness of the present series is a direct artefact of the sentiment--stance mismatch, and a stance-aware classifier would close it. The gated state-of-the-art models of Shah et al.\ \cite{shah2023trillion, shah2025wcb} are the most direct route, but a custom classifier fine-tuned on a manually annotated \gls{ecb} corpus is feasible at modest annotation cost. Second, large language models such as GPT-4o or Claude can serve as zero-shot stance annotators at near-human quality, providing a silver-standard label set for distillation into a smaller deployable encoder; reproducibility requires fixing temperature and logging the model version. Third, the daily \gls{ois} reaction used here can be replaced by the intraday windows of the Euro Area Monetary Policy Event-Study Database \cite{altavilla2019measuring}, separating the $14{:}15$~CET press-release window from the $14{:}45$--$15{:}15$~CET press-conference window; this directly tests the hypothesis that the \gls{qa} signal is the slower-incorporating of the two. Fourth, and most importantly from the policy perspective, the pipeline can be replicated on the press conferences of the Federal Reserve and the Bank of England under identical chunking, classification, and aggregation rules. This cross-bank replication is the planned next step in the collaboration with the Národná banka Slovenska and forms the empirical core of the envisaged co-authored publication: the cross-bank profile of communication asymmetry will directly inform the conditions under which a small open economy facing the euro area should treat \gls{ecb} sentiment as a leading indicator versus a contemporaneous signal.

The qualitative record of monetary policy can be read systematically, separated cleanly into its prepared and spontaneous components, and converted into a structured time series with measurable predictive value for the price of money. The framework developed in this thesis is one step in that programme; the limitations identified along the way describe extensions of scope rather than barriers to the central conclusion.